% Fragmento público generado desde propietarios íntegros.
% Residencia: 52.7.1.
% Fuente: ../incoming/radion_monografia_20260827/manuscrito/sections/08_complejo_electromagnetismo.tex:1-58
\noindent Les 3 régimes du TRIT admettent des réalisations complexe, électromagnétique et mécanique avec des types différenciés.\par\medskip

\subsubsection{Opérateur quadratique de racine interne de \(-1\)}

\paragraph{Représentation exponentielle du régime elliptique \(J^2=-I\)}

Avec \(J^2=-I\),
\begin{equation}
T_{\mathrm{ell}}(x,y)=e^{-xI-yJ}
=e^{-x}(\cos y\,I-\sin y\,J).
\label{eq:elliptic-exp}
\end{equation}
Sous l'identification générée par \(J\),
\[
T_{\mathrm{ell}}(x,y)\longleftrightarrow e^{-x-iy}.
\]
Le scalaire \(x\) contrôle la dilatation ; \(y\) contrôle la rotation orientée. La forme
\(a+bi\) se reconnaît comme élasticité plus rotation, non comme deux quantités
collées arbitrairement.

\paragraph{Représentation exponentielle du régime parabolique \(J^2=0\)}

Avec \(N^2=0\), l'exponentielle se tronque :
\begin{equation}
T_{\mathrm{par}}(x,y)=e^{-x}(I-yN).
\label{eq:parabolic-exp}
\end{equation}
Le régime décrit le seuil où il n'y a ni rotation périodique ni ouverture
exponentielle. Dans la lecture temporelle active, c'est le présent euclidien, mais
il conserve la mémoire des branches qui le délimitent.

\paragraph{Représentation exponentielle du régime hyperbolique \(J^2=+I\)}

Avec \(R^2=I\) et
\(P_\pm=(I\pm R)/2\),
\begin{equation}
T_{\mathrm{hyp}}(x,y)=e^{-xI-yR}
=q_+P_++q_-P_-,
\qquad q_\pm=e^{-x\mp y}.
\label{eq:hyperbolic-exp}
\end{equation}
Les invariants élémentaires sont
\begin{equation}
q_+q_-=e^{-2x},
\qquad
\frac{q_-}{q_+}=e^{2y}.
\label{eq:q-product-ratio}
\end{equation}
Le produit conserve le mode commun ; le quotient conserve la séparation
orientée.

Appliqué au couple angulaire,
\[
x=\frac{\pi A}{180},\qquad y=\frac{\pi C^*}{180}.
\]
C'est pourquoi \(A\) est la coordonnée commune de la publication et \(C^*\) son
contraste orienté. Ce n'est qu'après l'application d'un lecteur physique qu'ils sont reconnus
comme pôles électriques et magnétiques.

